\section{old section}

\subsection{Continuous analytical solutions in Fourier domain with true magnetization direction}
\label{subsec:Fourier_domain_true_mag}

Consider now that the 2D convolutions representing the 
$\alpha$-component $B_{\alpha}(\mathbf{r})$ (equation \ref{eq:B-alpha-Neumann}),
approximated total-field anomaly $B_{t}(\mathbf{r})$ (equation \ref{eq:B-t-Neumann}) 
and RTP anomaly $B_{r}(\mathbf{r})$ (equation \ref{eq:B-r-Neumann}) are 
evaluated at observation points $\mathbf{r}_{0} = \left[x \:\: y \:\: z_{0} \right]^{\top}$ 
located on a horizontal plane with constant vertical coordinate $z = z_{0}$. 
Then, we use the Fourier-Convolution Theorem \cite[][ p. 268]{blakely1996} to obtain:
\begin{equation}
\mathcal{F}\left[ B_{\alpha}(\mathbf{r}_{0}) \right] = 
\mathcal{F}\left[ \partial_{z} \Theta(\mathbf{r}'') \right] \: 
\mathcal{F}\left[ \partial_{\alpha h} \Psi(\mathbf{r}_{0}, \mathbf{r}'') \right] \: ,
\:\: \alpha = x, y, z \: ,
\label{eq:Fourier-B-alpha}
\end{equation}
\begin{equation}
\mathcal{F}\left[ B_{t}(\mathbf{r}_{0}) \right] = 
\mathcal{F}\left[ \partial_{z} \Theta(\mathbf{r}'') \right] \: 
\mathcal{F}\left[ \partial_{th} \Psi(\mathbf{r}_{0}, \mathbf{r}'') \right]
\label{eq:Fourier-B-t}
\end{equation}
and
\begin{equation}
\mathcal{F}\left[ B_{r}(\mathbf{r}_{0}) \right] = 
\mathcal{F}\left[ \partial_{z} \Theta(\mathbf{r}'') \right] \: 
\mathcal{F}\left[ \partial_{zz} \Psi(\mathbf{r}_{0}, \mathbf{r}'') \right] \: ,
\label{eq:Fourier-B-r}
\end{equation}
where $\mathcal{F}\left[ \cdot \right]$ denotes the 2D Fourier Transform (2DFT)
with respect to the horizontal coordinates $x$ and $y$ of the observation points.
The 2DFT's are two-dimensional functions of the wavenumbers $k_{x}$ and $k_{y}$ 
along the $x$ and $y$ directions, respectively. The wavenumbers have units of 
inverse distance. The dependence of the 2DFT's on $k_{x}$ and $k_{y}$ is conveniently
omitted to simplify the notation.

To define the 2DFT $\mathcal{F}\left[ \partial_{z} \Theta(\mathbf{r}'') \right]$ 
on the right side of equations \ref{eq:Fourier-B-t} and \ref{eq:Fourier-B-r}, 
where $\mathbf{r}'' = \left[x'' \:\: y'' \:\: z_{c} \right]^{\top}$ is the 
position vector defining a point $(x'', y'', z_{c})$ on the equivalent layer, 
we can first compute the derivative of $\Theta(\mathbf{r})$ (equation
\ref{eq:scalar-potential-Theta}) with respect to $z$ and then use 
\citet[][ p. 297]{blakely1996} to obtain:
\begin{equation}
\mathcal{F}\left[ \partial_{z} \Theta(\mathbf{r}'') \right] 
= 2\pi \, c_{m} \, e^{|k| \, z_{c}} 
\int\limits_{z_{c}}^{\infty} 
\mathcal{F}\left[ h(\mathbf{r}') \right] e^{-|k| \, z'} \:
dz' \: , \:\: z_{c} < z' \: ,
\label{eq:Fourier-Dz-Theta}
\end{equation}
where $|k| = \sqrt{k_{x}^{2} + k_{y}^{2}}$ and 
$\mathcal{F}\left[ h(\mathbf{r}') \right]$ represents the 
2DFT of the magnetization intensity distribution $h(\mathbf{r}')$
(equation \ref{eq:total-magnetization-vector}) on one horizontal slice
through the magnetic sources at $z = z'$.
In equation \ref{eq:Fourier-Dz-Theta}, we assume that 
$h(\mathbf{r}')$ is zero outside the volume $v$ of the sources
and, consequently, at any coordinate $z < z_{0} < z_{c}$.
According to \citet[][ p. 278]{blakely1996}, the 2DFT's 
$\mathcal{F}\left[ \partial_{\alpha h} \Psi(\mathbf{r}_{0}, \mathbf{r}'') \right]$ 
(equation \ref{eq:Fourier-B-alpha}),
$\mathcal{F}\left[ \partial_{th} \Psi(\mathbf{r}_{0}, \mathbf{r}'') \right]$ 
(equation \ref{eq:Fourier-B-t}) and 
$\mathcal{F}\left[ \partial_{zz} \Psi(\mathbf{r}_{0}, \mathbf{r}'') \right]$ 
(equation \ref{eq:Fourier-B-r}) are given by:
\begin{equation}
\mathcal{F}\left[ \partial_{\alpha h} \Psi(\mathbf{r}_{0}, \mathbf{r}'') \right] = 
\mathcal{U}_{\alpha} \: \mathcal{U}_{h} \: |k| \, e^{|k| \, (z_{0} - z_{c})}
\: , \:\: z_{0} < z_{c} \: , \:\: \alpha = x, y, z \: ,
\label{eq:Fourier-Dalphah-Psi}
\end{equation}
\begin{equation}
\mathcal{F}\left[ \partial_{th} \Psi(\mathbf{r}_{0}, \mathbf{r}'') \right] = 
\mathcal{U}_{t} \: \mathcal{U}_{h} \: |k| \, e^{|k| \, (z_{0} - z_{c})}
\: , \:\: z_{0} < z_{c} \: ,
\label{eq:Fourier-Dth-Psi}
\end{equation}
and
\begin{equation}
\mathcal{F}\left[ \partial_{zz} \Psi(\mathbf{r}_{0}, \mathbf{r}'') \right] 
= |k| \, e^{|k| \, (z_{0} - z_{c})} \: ,
\:\: z_{0} < z_{c} \: ,
\label{eq:Fourier-Dzz-Psi}
\end{equation}
where 
$\mathcal{U}_{x} \equiv \mathcal{U}(0^{\circ}, 0^{\circ})$,
$\mathcal{U}_{y} \equiv \mathcal{U}(0^{\circ}, 90^{\circ})$,
$\mathcal{U}_{z} \equiv \mathcal{U}(90^{\circ}, 0^{\circ})$,
$\mathcal{U}_{t} \equiv \mathcal{U}(I_{t}, D_{t})$ and
$\mathcal{U}_{h} \equiv \mathcal{U}(I_{h}, D_{h})$ are defined by the complex function
\begin{equation}
\mathcal{U}(I, D) = \hat{u}_{z}(I, D) + 
\mathrm{i} \: \frac{ \hat{u}_{x}(I, D) \: k_{x} + 
\hat{u}_{y}(I, D) \: k_{y} }{|k|} \: ,
\label{eq:U-complex-function}
\end{equation}
in terms of the imaginary unity $\mathrm{i} = \sqrt{-1}$
and the Cartesian components of the unit vector
$\hat{\mathbf{u}}(I, D)$ (equation \ref{eq:u-unit-vector}) computed with 
the corresponding inclination $I$ and declination $D$ values.
Now, by substituting the 2DFT's 
$\mathcal{F}\left[ \partial_{z} \Theta(\mathbf{r}'') \right]$ 
(equation \ref{eq:Fourier-Dz-Theta}),
$\mathcal{F}\left[ \partial_{\alpha h} \Psi(\mathbf{r}_{0}, \mathbf{r}'') \right]$
(equation \ref{eq:Fourier-Dalphah-Psi}), 
$\mathcal{F}\left[ \partial_{th} \Psi(\mathbf{r}_{0}, \mathbf{r}'') \right]$
(equation \ref{eq:Fourier-Dth-Psi}) and
$\mathcal{F}\left[ \partial_{zz} \Psi(\mathbf{r}_{0}, \mathbf{r}'') \right]$ 
(equation \ref{eq:Fourier-Dzz-Psi}) into those defining
$\mathcal{F}\left[ B_{t}(\mathbf{r}_{0}) \right]$
(equation \ref{eq:Fourier-B-alpha}),
$\mathcal{F}\left[ B_{t}(\mathbf{r}_{0}) \right]$
(equation \ref{eq:Fourier-B-t}) and 
$\mathcal{F}\left[ B_{r}(\mathbf{r}_{0}) \right]$
(equation \ref{eq:Fourier-B-r}), we obtain:
\begin{equation}
\mathcal{F}\left[ B_{\alpha}(\mathbf{r}_{0}) \right] = 
2\pi \, c_{m} \, 
\mathcal{U}_{\alpha} \: \mathcal{U}_{h} \: |k| \, e^{|k| \, z_{0}} 
\int\limits_{z_{0}}^{\infty} 
\mathcal{F}\left[ h(\mathbf{r}') \right] e^{-|k| \, z'} \:
dz' \: , \:\: z_{0} < z' \: , \:\: \alpha = x, y, z \: ,
\label{eq:Fourier-B-alpha-final}
\end{equation}
\begin{equation}
\mathcal{F}\left[ B_{t}(\mathbf{r}_{0}) \right] = 
2\pi \, c_{m} \, 
\mathcal{U}_{t} \: \mathcal{U}_{h} \: |k| \, e^{|k| \, z_{0}} 
\int\limits_{z_{0}}^{\infty} 
\mathcal{F}\left[ h(\mathbf{r}') \right] e^{-|k| \, z'} \:
dz' \: , \:\: z_{0} < z' \: ,
\label{eq:Fourier-B-t-final}
\end{equation}
and
\begin{equation}
	\mathcal{F}\left[ B_{r}(\mathbf{r}_{0}) \right] = 
	2\pi \, c_{m}
	\: |k| \, e^{|k| \, z_{0}} 
	\int\limits_{z_{0}}^{\infty} 
	\mathcal{F}\left[ h(\mathbf{r}') \right] e^{-|k| \, z'} \:
	dz' \: , \:\: z_{0} < z' \: .
	\label{eq:Fourier-B-r-final}
\end{equation}
Note that, in deducing equations 
\ref{eq:Fourier-B-alpha-final}--\ref{eq:Fourier-B-r-final}, 
we have replaced $z_{c}$ by
$z_{0}$ in the 2DFT $\mathcal{F}\left[ \partial_{z} \Theta(\mathbf{r}'') \right]$ 
(equation \ref{eq:Fourier-Dz-Theta}). This change does not affect the result
because we assume that $h(\mathbf{r}')$ is zero outside the volume $v$ of the sources
and, consequently, at any coordinate $z < z_{0} < z_{c}$.

Regardless of notation differences, the 2DFT's
$\mathcal{F}\left[ B_{\alpha}(\mathbf{r}_{0}) \right]$ 
(equation \ref{eq:Fourier-B-alpha-final}) and 
$\mathcal{F}\left[ B_{t}(\mathbf{r}_{0}) \right]$ (equation \ref{eq:Fourier-B-t-final})
are exactly the same found, for example, in \citet{gunn1975} or 
\citet[][ p. 298]{blakely1996}.
It is important to stress, however, that differently from those authors, 
we have not deduced equations \ref{eq:Fourier-B-alpha-final} 
and \ref{eq:Fourier-B-t-final} by computing the 2DFT's of the volume integrals 
$B_{\alpha}(\mathbf{r})$ (equation \ref{eq:B-alpha}) and 
$B_{t}(\mathbf{r})$ (equation \ref{eq:B-t}) at the horizontal plane $z = z_{0}$. 
Instead, we have deduced them by computing the 2DFT's of the surface integrals
given by equations \ref{eq:B-alpha-Neumann} and \ref{eq:B-t-Neumann}.
It can also be verified from equations 
\ref{eq:Fourier-Dz-Theta}, \ref{eq:Fourier-B-t-final} and 
\ref{eq:Fourier-B-r-final} that $\mathcal{F}\left[ B_{r}(\mathbf{r}_{0}) \right]$ 
(equation \ref{eq:Fourier-B-r-final}) can be rewritten as follows:
\begin{equation}
\mathcal{F}\left[ B_{r}(\mathbf{r}_{0}) \right]
= \frac{1}{\mathcal{U}_{t} \,\mathcal{U}_{h}} \, 
\mathcal{F}\left[ B_{t}(\mathbf{r}_{0}) \right] \: ,
\label{eq:Fourier-RTP}
\end{equation}
which represents the common expression defining the RTP in Fourier domain
\cite[e.g.,][ p. 331]{blakely1996}.

The results presented here establish, for the first time, 
the theoretical connection between the RTP in space and Fourier domains.
We show that the well-known equation \ref{eq:Fourier-RTP} for computing the
RTP in Fourier domain can be obtained from the 2DFT's of the surface integrals
defining the RTP in space domain (equations \ref{eq:B-t-Neumann} and \ref{eq:B-r-Neumann}).
We stress that these equations are valid for the particular case in which the
equivalent layer of dipoles has the same uniform magnetization direction
$\hat{\mathbf{u}}(I_{h}, D_{h})$ of the true sources 
(equation \ref{eq:total-magnetization-vector}) and the observation points 
$\mathbf{r}_{0} = \left[x \:\: y \:\: z_{0} \right]^{\top}$ 
located on a horizontal plane with constant vertical coordinate $z = z_{0}$.


\subsection{Continuous analytical solutions with arbitrary uniform magnetization direction}
\label{subsec:arbitrary_mag}


Alternative analytical expressions for the 2DFT's 
$\mathcal{F}\left[ B_{\alpha}(\mathbf{r}_{0}) \right]$ 
(equations \ref{eq:Fourier-B-alpha} and \ref{eq:Fourier-B-alpha-final})
and
$\mathcal{F}\left[ B_{t}(\mathbf{r}_{0}) \right]$ 
(equations \ref{eq:Fourier-B-t} and \ref{eq:Fourier-B-t-final}) 
can also be deduced for the case in which the
equivalent layer has an arbitrary uniform magnetization direction 
$\hat{\mathbf{u}}(I_{p}, D_{p})$. These alternative expressions are obtained from the
simple manipulation of equations \ref{eq:Fourier-Dz-Theta},
\ref{eq:Fourier-Dalphah-Psi} and \ref{eq:Fourier-Dth-Psi}, as follows:
\begin{equation}
\mathcal{F}\left[ B_{\alpha}(\mathbf{r}_{0}) \right] = 
\mathcal{F}\left[ p(\mathbf{r}'') \right] \: 
\mathcal{F}\left[ \partial_{\alpha p} \Psi(\mathbf{r}_{0}, \mathbf{r}'') \right] \: ,
\:\: \alpha = x, y, z \: ,
\label{eq:Fourier-B-alpha-arbitrary}
\end{equation}
and
\begin{equation}
\mathcal{F}\left[ B_{t}(\mathbf{r}_{0}) \right] = 
\mathcal{F}\left[ p(\mathbf{r}'') \right] \: 
\mathcal{F}\left[ \partial_{tp} \Psi(\mathbf{r}_{0}, \mathbf{r}'') \right] \: ,
\label{eq:Fourier-B-t-arbitrary}
\end{equation}
where
\begin{equation}
\mathcal{F}\left[ p(\mathbf{r}'') \right] 
= \frac{\mathcal{U}_{h}}{\mathcal{U}_{p}} \, 
\mathcal{F}\left[ \partial_{z} \Theta(\mathbf{r}'') \right] \: ,
\label{eq:Fourier-p}
\end{equation}
\begin{equation}
\mathcal{F}\left[ \partial_{\alpha p} \Psi(\mathbf{r}_{0}, \mathbf{r}'') \right] = 
\mathcal{U}_{\alpha} \: \mathcal{U}_{p} \: |k| \, e^{|k| \, (z_{0} - z_{c})}
\: , \:\: z_{0} < z_{c} \: , \:\: \alpha = x, y, z \: ,
\label{eq:Fourier-Dalphap-Psi}
\end{equation}
and
\begin{equation}
\mathcal{F}\left[ \partial_{tp} \Psi(\mathbf{r}_{0}, \mathbf{r}'') \right] = 
\mathcal{U}_{t} \: \mathcal{U}_{p} \: |k| \, e^{|k| \, (z_{0} - z_{c})}
\: , \:\: z_{0} < z_{c} \: ,
\label{eq:Fourier-Dtp-Psi}
\end{equation}
with $\mathcal{U}_{p} \equiv \mathcal{U}(I_{p}, D_{p})$ (equation 
\ref{eq:U-complex-function}) defined in terms of the arbitrary constant 
inclination $I_{p}$ and declination $D_{p}$ values.
Note that, in this case, it is not possible to compute the 2DFT 
$\mathcal{F}\left[ B_{r}(\mathbf{r}_{0}) \right]$ of the RTP anomaly from the
2DFT $\mathcal{F}\left[ p(\mathbf{r}'') \right]$ of the generic physical-property 
distribution $p(\mathbf{r}'')$ because, according to equations \ref{eq:Fourier-B-r}, 
\ref{eq:Fourier-Dzz-Psi} and \ref{eq:Fourier-p}:
\begin{equation}
\mathcal{F}\left[ B_{r}(\mathbf{r}_{0}) \right]
\ne \mathcal{F}\left[ p(\mathbf{r}'') \right] \, 
\mathcal{F}\left[ \partial_{zz} \Psi(\mathbf{r}_{0}, \mathbf{r}'') \right] \: .
\label{eq:Fourier-B-r-wrong}
\end{equation}

The alternative analytical expressions presented here 
(equations \ref{eq:Fourier-B-alpha-arbitrary}--\ref{eq:Fourier-B-r-wrong})
show, in Fourier domain, that the planar equivalent layer with 
generic magnetic moment distribution 
$p(\mathbf{r}'')$ obtained by using an arbitrary uniform magnetization direction
$\hat{\mathbf{u}}(I_{p}, D_{p})$ can exactly reproduce the magnetic field components
$B_{\alpha}(\mathbf{r}_{0})$ and approximated total-field anomaly
$B_{t}(\mathbf{r}_{0})$, but cannot be used to compute the RTP anomaly
$B_{r}(\mathbf{r}_{0})$.
